\documentclass[../RFC-HDFG-2026-003.tex]{subfiles}

\begin{document}

\section{Implementation Notes}
\label{sec:impl}

This section maps the design in the preceding sections to the library
source files it touches.

\begin{center}
\small
\renewcommand{\arraystretch}{1.25}
\begin{tabularx}{\textwidth}{lX}
\toprule
\textbf{File} & \textbf{Change} \\
\midrule
\texttt{src/H5Zdevelop.h} & Add \texttt{H5Z\_init\_func\_t} and
  \texttt{H5Z\_term\_func\_t}, and append \texttt{init} and \texttt{term}
  to \texttt{H5Z\_class3\_t} after \texttt{description}
  (\SectionRef{sec:filter-state}); add \texttt{H5Z\_blob\_loc\_t} and the
  four blob callback typedefs, and append the four fields after
  \texttt{term} (\SectionRef{sec:blob-callbacks}). \\
\texttt{src/H5Zprivate.h} & Add \texttt{state} and \texttt{state\_status}
  (none / ready / failed) to \texttt{H5Z\_filter\_info\_t}, so the state
  travels with the dataset's own pipeline and \texttt{H5Z\_pipeline}
  reaches it at every one of its call sites without new arguments
  (\SectionRef{sec:filter-state}).  Add the reference-counted
  \texttt{H5Z\_blob\_buf\_t} wrapper and the \texttt{aux}/\texttt{aux\_loc}
  fields to \texttt{H5Z\_filter\_info\_t}
  (\SectionRef{sec:blob-internal-state}). \\
\texttt{src/H5Z.c} & \texttt{H5Z\_state\_init} (runs \texttt{init} per
  entry whose class is available; strict at create and copy, lenient at
  open) and \texttt{H5Z\_state\_term}; \texttt{H5Z\_pipeline} passes the
  entry's state, runs a deferred \texttt{init} for an entry whose plugin
  it has just loaded, and refuses to run an entry whose class has
  \texttt{init} but whose state is not ready; \texttt{H5Z\_can\_apply\_direct}
  refuses such classes for fractal heaps (\SectionRef{sec:filter-state}).
  Wire up \texttt{write\_blob}/\texttt{read\_blob}/
  \texttt{delete\_blob}/\texttt{close\_blob} (or the default
  \texttt{H5HG}-based implementation) into the create/open/close/delete
  paths (\SectionRef{sec:blob-lifecycle}), reading and writing blobs
  before \texttt{init} runs at open; reject partial storage
  callback sets in \texttt{H5Z\_register3}; never read a custom locator
  through \texttt{H5HG\_read} when the plugin is unavailable;
  segmented default storage for blobs over 1~MiB
  (\SectionRef{sec:blob-storage-format}); \texttt{lookup3} checksum on
  write and verification on read; \texttt{H5Z\_blob\_buf\_new}/
  \texttt{\_incref} and the decrement-and-maybe-free logic in
  \texttt{H5Z\_blob\_release}; prototype and validate the collective
  identical-insert path for parallel jobs, including a debug-build
  locator-consistency check (\SectionRef{sec:blob-parallel}). \\
\texttt{src/H5Fprivate.h}, \texttt{src/H5Fint.c} & A per-open-file table
  of stored default-storage blobs (size, checksum, buffer reference,
  locator) for content sharing, released at file close
  (\SectionRef{sec:blob-sharing}). \\
\texttt{src/H5Pocpl.c} & Implement \texttt{H5Pappend\_filter\_blob}
  (\SectionRef{sec:append-filter-blob}); extend
  \texttt{H5P\_\_ocrt\_pipeline\_enc}/\texttt{\_dec} to serialize the
  \texttt{aux} buffer inline (\SectionRef{sec:blob-encode-decode}). \\
\texttt{src/H5Opline.c} & Pipeline-message encode/decode of the
  \texttt{H5O\_PLINE\_EXT\_BLOB} extension block carrying
  \texttt{aux\_loc}, registered in the block-type table
  RFC-HDFG-2026-001 defines (\SectionRef{sec:blob-format}), including
  the storage-kind byte, checksum, and size; the \texttt{delete} handler
  decrements the blob object's reference count and removes it (every
  segment, for a segmented blob) at zero, or calls \texttt{delete\_blob}
  for custom storage (\SectionRef{sec:blob-sharing});
  \texttt{H5O\_\_pline\_copy} shares the refcounted \texttt{aux} buffer
  (\texttt{H5Z\_blob\_buf\_incref}) instead of deep-copying it
  (\SectionRef{sec:blob-repack}) and clears the state fields, so property
  lists and other copies never carry state (\SectionRef{sec:filter-state}). \\
\texttt{src/H5Dint.c} & \texttt{init} in \texttt{H5D\_\_create} before
  \texttt{H5D\_\_update\_oh\_info}, whose early allocation filters fill
  chunks; \texttt{init} in \texttt{H5D\_\_open\_oid} after the pipeline and
  any blob are read; \texttt{term} in \texttt{H5D\_close} after the
  layout's chunk-cache flush, and on the create/open failure paths
  (\SectionRef{sec:filter-state}). \\
\texttt{src/H5Dchunk.c} & \texttt{H5D\_\_chunk\_copy} brackets the
  variable-length/reference copy path with \texttt{init}/\texttt{term} on
  a private copy of the pipeline (\SectionRef{sec:filter-state}).  No
  blob-related changes: the blob is configuration data loaded once at open
  time. \\
\texttt{tools/src/h5repack/} & No changes needed; see
  \SectionRef{sec:blob-repack}. \\
\texttt{test/} & New tests: per-dataset state lifecycle
  (\SectionRef{sec:test-filter-state}), create/open/close, parallel write/broadcast,
  \texttt{H5Pencode}/\texttt{H5Pdecode} round-trip, and \texttt{h5repack}
  cross-file copy for a blob-configured filter (\SectionRef{sec:test-blob}). \\
\bottomrule
\end{tabularx}
\end{center}

\end{document}
